% Unit 057 terjemahan Bahasa Indonesia dari chapter4.tex baris 1513--1680.
% Sumber: Methods of Algebra, Volume 2, commit
% 9a5803ff77dd3257484cb177f851a73770a59dd3 (CC BY 4.0).
% stable-unit-id: o014.aljabr2.chapter4.derived-functors-general
% Cakupan: Bagian 4.7 lengkap; berhenti sebelum Bagian 4.8 pada baris 1681.

% segment-id: o014.li2.chapter4.derived-functors-general.h001
\section{Teori Umum Funktor Turunan}\label{sec:derived-functors}
\hypertarget{o014.aljabr2.chapter4.derived-functors-general}{ }

% segment-id: o014.li2.chapter4.derived-functors-general.p001
Misalkan $\mathcal{A}$ dan $\mathcal{A}'$ kategori abelian, dan tetapkan
funktor aditif $F:\mathcal{A}\to\mathcal{A}'$. Sekarang kita menerapkan teori
dalam \S\ref{sec:triangulated-functor-localization}, khususnya situasi
\eqref{eqn:derived-situation}, pada
% segment-id: o014.li2.chapter4.derived-functors-general.g001
\begin{equation*}\begin{tikzcd}
	\cate{K}^{\star}(\mathcal{A}) \arrow[r, "{\cate{K}^\star F}"] \arrow[d, "Q"'] & \cate{K}^{\star}(\mathcal{A}') \arrow[d, "{Q'}"] \\
	\cate{D}^{\star}(\mathcal{A}) & \cate{D}^{\star}(\mathcal{A}')
\end{tikzcd} \qquad (\star \in \{+, -, \hspace{0.8em}\}) \end{equation*}
Di sini $Q$ dan $Q'$ adalah funktor lokalisasi kategori bertriangulasi, dan
$\cate{K}^\star F$ merupakan funktor bertriangulasi
(Proposisi~\ref{prop:functor-induces-triangulated}). Superskrip kosong
$\star=\hspace{0.8em}$ bersesuaian dengan kasus $\cate{K}(\mathcal{A})$ dan
$\cate{D}(\mathcal{A})$. Ingat bahwa
$\cate{D}^{\star}(\mathcal{A})=\cate{K}^{\star}(\mathcal{A})/\mathcal{N}^{\star}$,
dengan $\mathcal{N}^{\star}$ subkategori bertriangulasi jenuh yang terdiri atas
kompleks asiklik; demikian pula untuk $\mathcal{A}'$.

% segment-id: o014.li2.chapter4.derived-functors-general.p002
Kecuali jika $F$ eksak, secara umum tidak ada funktor bertriangulasi
$\cate{D}^{\star}(\mathcal{A})\to\cate{D}^{\star}(\mathcal{A}')$ yang membuat
diagram di atas komutatif. Pendekatan yang wajar adalah mengkaji funktor
turunan kiri dan kanan dalam pengertian
Definisi~\ref{def:Verdier-localization-functor}, masing-masing melalui dua
jenis ekstensi Kan.

% segment-id: o014.li2.chapter4.derived-functors-general.d001
\begin{definisi}[Funktor turunan antarkategori turunan]\label{def:derived-functor}
	\index{funktor turunan}
	\index[sym1]{RF@$\mathrm{R}F$, ${}^\star \mathrm{R}F$}
	\index[sym1]{LF@$\mathrm{L}F$, ${}^\star \mathrm{L}F$}
	\index[sym1]{RnF} \index[sym1]{LnF}
	Untuk $\star\in\{+,-,\hspace{0.8em}\}$, jika
	\begin{align*}
		{}^\star \mathrm{R}F: \cate{D}^\star(\mathcal{A}) & \to \cate{D}^\star(\mathcal{A}'): \quad \text{funktor turunan kanan dari $\cate{K}^\star F$}, \quad \text{atau} \\
		{}^\star \mathrm{L}F: \cate{D}^\star(\mathcal{A}) & \to \cate{D}^\star(\mathcal{A}'): \quad \text{funktor turunan kiri dari $\cate{K}^\star F$}
	\end{align*}
	ada, maka funktor tersebut disebut \emph{funktor turunan kanan} (atau
	\emph{funktor turunan kiri}) dari $F:\mathcal{A}\to\mathcal{A}'$. Jika
	funktor turunannya ada, untuk setiap $n\in\Z$ didefinisikan
	\begin{gather*}
		{}^\star \mathrm{R}^n F := \Hm^n \circ {}^\star \mathrm{R}F, \\
		{}^\star \mathrm{L}_n F := \Hm^{-n} \circ {}^\star \mathrm{L}F.
	\end{gather*}
	Keduanya merupakan funktor bernilai di $\mathcal{A}'$. Subskrip $n$
	digunakan untuk menyesuaikan dengan notasi kompleks rantai yang lazim ketika
	mengkaji funktor turunan kiri. Superskrip kiri $\star$ sering dihilangkan
	jika tidak menimbulkan kerancuan.
\end{definisi}

% segment-id: o014.li2.chapter4.derived-functors-general.e001
\begin{contoh}[Menurunkan funktor eksak]
	Kasus paling sederhana adalah ketika $F$ merupakan funktor eksak. Untuk
	setiap $\star\in\{+,-,\hspace{0.8em}\}$, funktor turunan
	${}^\star\mathrm{R}F$ dan ${}^\star\mathrm{L}F$ selalu ada. Memang,
	$\cate{K}^\star F$ mempertahankan kompleks asiklik, sehingga funktor
	turunannya langsung ditentukan oleh sifat universal lokalisasi.
\end{contoh}

% segment-id: o014.li2.chapter4.derived-functors-general.t001
\begin{teorema}[Barisan eksak panjang funktor turunan]\label{prop:derived-long-exact}
	\index{barisan eksak panjang}
	Tetapkan $\star\in\{+,-,\hspace{0.8em}\}$. Jika
	${}^\star\mathrm{R}F$ atau ${}^\star\mathrm{L}F$ ada, maka setiap segitiga
	terbedakan $X\to Y\to Z\xrightarrow{+1}$ dalam
	$\cate{D}^\star(\mathcal{A})$ menghasilkan barisan eksak panjang kanonik
	\begin{gather*}
		\cdots \to {}^\star \mathrm{R}^{n-1} F(Z) \to {}^\star \mathrm{R}^n F(X) \to {}^\star \mathrm{R}^n F(Y) \to {}^\star \mathrm{R}^n F(Z) \to {}^\star \mathrm{R}^{n+1} F(X) \to \cdots , \\
		\cdots \to {}^\star \mathrm{L}_{n+1} F(Z) \to {}^\star \mathrm{L}_n F(X) \to {}^\star \mathrm{L}_n F(Y) \to {}^\star \mathrm{L}_n F(Z) \to {}^\star \mathrm{L}_{n-1} F(X) \to \cdots .
	\end{gather*}
\end{teorema}
% segment-id: o014.li2.chapter4.derived-functors-general.pf001
\begin{bukti}
	Menurut definisi, ${}^\star\mathrm{R}F$ maupun ${}^\star\mathrm{L}F$
	merupakan funktor bertriangulasi. Selanjutnya, terapkan funktor kohomologi
	$\Hm^0$.
\end{bukti}

% segment-id: o014.li2.chapter4.derived-functors-general.p003
Bagaimana kita menjamin keberadaan funktor turunan dan menghitungnya?
Proposisi~\ref{prop:localization-injective} sudah memuat gagasan umumnya.
Untuk $\star\in\{+,-,\hspace{0.8em}\}$, ambil $\mathcal{N}^\star$ seperti
dalam Definisi~\ref{def:derived-cat}. Kita mencari subkategori bertriangulasi
$\mathcal{I}$ dari $\cate{K}^\star(\mathcal{A})$ agar ${}^\star\mathrm{R}F$
atau ${}^\star\mathrm{L}F$ dapat didefinisikan sebagai komposisi pada diagram
funktor
% segment-id: o014.li2.chapter4.derived-functors-general.g002
\[ \cate{D}^\star(\mathcal{A}) = \cate{K}^\star(\mathcal{A})/\mathcal{N}^\star \xleftarrow{\text{ekuivalensi}} \mathcal{I}/\mathcal{I} \cap \mathcal{N}^\star \xrightarrow{\text{sifat universal lokalisasi}} \cate{K}^\star(\mathcal{A}')/\mathcal{N}^\star = \cate{D}^\star(\mathcal{A}') . \]
Bagaimana $\mathcal{I}$ harus dipilih? Subkategori tersebut hendaknya
$F$-injektif (atau $F$-projektif) dalam pengertian berikut.

% segment-id: o014.li2.chapter4.derived-functors-general.c001
\begin{konvensi}
	\index{subkategori injektif-$\cdots$}
	\index{subkategori projektif-$\cdots$}
	Tetapkan $\star$. Jika subkategori bertriangulasi $\mathcal{I}$ dari
	$\cate{K}^\star(\mathcal{A})$ merupakan subkategori bertriangulasi
	$\cate{K}^\star F$-injektif (atau $\cate{K}^\star F$-projektif) dalam
	pengertian Definisi~\ref{def:F-injective}, maka untuk singkatnya
	$\mathcal{I}$ disebut subkategori $F$-injektif (atau $F$-projektif).
\end{konvensi}

% segment-id: o014.li2.chapter4.derived-functors-general.p004
Jika terdapat subkategori $F$-injektif (atau $F$-projektif) $\mathcal{I}$
dari $\cate{K}^\star(\mathcal{A})$,
Proposisi~\ref{prop:localization-injective} menjamin bahwa
${}^\star\mathrm{R}F$ (atau ${}^\star\mathrm{L}F$) ada. Pembatasannya ke
$\mathcal{I}$ isomorfik dengan $Q'\cate{K}^\star F$, sedangkan rumus limit
dalam Catatan~\ref{rem:Deligne-derived-functor} memperjelas morfisme kanonik
yang dibawa funktor turunan sebagai ekstensi Kan:
% segment-id: o014.li2.chapter4.derived-functors-general.g003
\begin{equation}\label{eqn:L-R-can}
	Q' \cate{K}^\star F \to ({}^\star \mathrm{R}F)Q \quad \text{atau} \quad ({}^\star \mathrm{L}F)Q \to Q' (\cate{K}^\star F).
\end{equation}

% segment-id: o014.li2.chapter4.derived-functors-general.p005
Seluruh teori umum yang dibangun dalam
\S\ref{sec:triangulated-functor-localization} dapat langsung digunakan.
Sebagai contoh, tinjau funktor aditif antarkategori abelian
\[ \mathcal{A} \xrightarrow{F} \mathcal{A}' \xrightarrow{F'} \mathcal{A}'' . \]
Jika funktor turunan yang dibahas ada, terdapat morfisme kanonik
\[ {}^\star \mathrm{R}(F'F) \to ({}^* \mathrm{R}F')({}^* \mathrm{R}F), \quad \left( {}^\star \mathrm{L}F' \right) \left( {}^\star \mathrm{L}F \right) \to {}^\star \mathrm{L}(F' F), \]
dan Teorema~\ref{prop:localization-triangulated-composite} memberikan syarat
cukup agar morfisme tersebut merupakan isomorfisme.

% segment-id: o014.li2.chapter4.derived-functors-general.p006
Selain itu, terdapat kompatibilitas berikut.

% segment-id: o014.li2.chapter4.derived-functors-general.l001
\begin{lema}\label{prop:derived-functor-res}
	Jika $\cate{K}(\mathcal{A})$ dan $\cate{K}^+(\mathcal{A})$ masing-masing
	mempunyai subkategori $F$-injektif, maka pembatasan $\mathrm{R}F$ ke
	$\cate{D}^+(\mathcal{A})$ memberikan ${}^+\mathrm{R}F$.

	Pernyataan yang bersesuaian juga berlaku bagi subkategori $F$-projektif,
	$\mathrm{L}F$, dan ${}^-\mathrm{L}F$.
\end{lema}
% segment-id: o014.li2.chapter4.derived-functors-general.pf002
\begin{bukti}
	Kita terlebih dahulu membahas $\mathrm{R}F$. Misalkan
	$X\in\Obj(\cate{K}^+(\mathcal{A}))$. Tuliskan $\mathrm{qis}_{X/}$ untuk
	kategori semua kuasi-isomorfisme $X\to Y$ dalam $\cate{K}(\mathcal{A})$,
dan $\mathrm{qis}^+_{X/}$ untuk versinya di $\cate{K}^+(\mathcal{A})$.
Catatan~\ref{rem:Deligne-derived-functor} memberikan isomorfisme kanonik
	\begin{equation}\label{eqn:derived-as-limit}\begin{aligned}
		\mathrm{R}F(QX) \simeq \varinjlim_{[X \to Y] \in \Obj(\mathrm{qis}_{X/})} (Q' \cate{K}F)(Y).
	\end{aligned}\end{equation}

	Untuk setiap objek $X\to Y$ dalam $\mathrm{qis}_{X/}$, jika $m\ll0$
	(hanya bergantung pada $X$), maka
	$X\to Y\to\tau^{\geq m}Y$ merupakan objek dari
	$\mathrm{qis}^+_{X/}$. Lema~\ref{prop:S-filtered} menjamin bahwa
	$\mathrm{qis}_{X/}$ tersaring, sehingga
	Proposisi~\ref{prop:cofinal-filtered} menjamin bahwa subkategori penuh
	$\mathrm{qis}^+_{X/}$ kofinal di dalamnya. Oleh karena itu, $\varinjlim$
dalam \eqref{eqn:derived-as-limit} dapat diambil di atas
$\mathrm{qis}^+_{X/}$ (Proposisi~\ref{prop:cofinal-lim}), dan hasilnya tepat
${}^+\mathrm{R}F(QX)$.

	Untuk $\mathrm{L}F$, gunakan versi funktor turunan kiri dari
	Catatan~\ref{rem:Deligne-derived-functor} dan lakukan pemenggalan dalam arah
	sebaliknya.
\end{bukti}

% segment-id: o014.li2.chapter4.derived-functors-general.p007
Selanjutnya kita membahas bifunktor yang diperkenalkan dalam
Konvensi~\ref{con:bifunctor}. Teori yang sesuai juga telah disiapkan dalam
\S\ref{sec:triangulated-functor-localization}. Misalkan $\mathcal{A}_1$,
$\mathcal{A}_2$, dan $\mathcal{B}$ kategori abelian, dan bifunktor
$F:\mathcal{A}_1\times\mathcal{A}_2\to\mathcal{B}$ aditif dalam setiap
variabel. Seperti dalam Definisi--Proposisi~\ref{def:bifunctor-cplx}, definisikan
% segment-id: o014.li2.chapter4.derived-functors-general.g004
\begin{equation}\label{eqn:totCF}\begin{aligned}
		\cate{C}_{\oplus} F & := \tot_{\oplus} \circ \cate{C}^2 F \quad \text{(jika $\mathcal{B}$ mempunyai koproduk terhitung)}, \\
		\cate{C}_{\Pi} F & := \tot_{\Pi} \circ \cate{C}^2 F \quad \text{(jika $\mathcal{B}$ mempunyai produk terhitung)}.
\end{aligned}\end{equation}
Menurut Proposisi~\ref{prop:bifunctor-cplx-homotopy}, keduanya mempunyai
faktorisasi melalui funktor pada tingkat $\cate{K}(\cdot)$, yaitu
$\cate{K}_{\oplus}F$ dan $\cate{K}_{\Pi}F$.

% Reference: [KS06, Proposition 11.2.11]
% segment-id: o014.li2.chapter4.derived-functors-general.pr001
\begin{proposisi}
	Funktor
	$\cate{K}_{\Pi}F,\;\cate{K}_{\oplus}F:
	\cate{K}^\star(\mathcal{A}_1)\times\cate{K}^\star(\mathcal{A}_2)
	\to\cate{K}^\star(\mathcal{B})$ di atas merupakan bifunktor bertriangulasi
dalam pengertian Definisi~\ref{def:triangulated-bifunctor}.
\end{proposisi}
% segment-id: o014.li2.chapter4.derived-functors-general.pf003
\begin{bukti}
	Ambil $\cate{K}_{\Pi}F$ sebagai contoh (dengan asumsi $\mathcal{B}$
	mempunyai produk terhitung). Pembahasan setelah
	Proposisi~\ref{prop:bifunctor-cplx-homotopy} menunjukkan bahwa funktor ini
	kompatibel dengan pergeseran dalam setiap variabel, sedangkan
	Proposisi~\ref{prop:tot-shift} memberikan diagram antikomutatif yang
	dipersyaratkan dalam definisi bifunktor bertriangulasi. Tinggal dibuktikan
	bahwa $\cate{C}_{\Pi}F$ mempertahankan kerucut pemetaan dalam setiap
	variabel. Untuk variabel pertama, ambil morfisme
	$f:X_1\to X'_1$ dalam $\cate{C}(\mathcal{A}_1)$ dan objek $X_2$ dari
	$\cate{C}(\mathcal{A}_2)$. Terdapat isomorfisme kanonik dalam $\mathcal{B}$
	\begin{align*}
		\cate{C}_{\Pi}F(\Cone(f), X_2)^n & \simeq \prod_{p+q=n} F(X_1^{p+1}, X_2^q ) \times F( (X'_1)^p, X_2^q ) \\
		& \simeq \prod_{p+q=n+1} F( X_1^p, X_2^q ) \times \prod_{p+q=n} F( (X'_1)^p, X_2^q ) \\
		& \simeq \Cone\left(\cate{C}_{\Pi}(f, \identity_{X_2}) \right)^n , \quad n \in \Z.
	\end{align*}
	Tinggal diverifikasi bahwa isomorfisme-isomorfisme tersebut
	membentuk isomorfisme kompleks dan kompatibel dengan morfisme
	$\alpha(\cdot)$ dan $\beta(\cdot)$ yang menyertai kerucut pemetaan. Tidak
	ada kesulitan mendasar; perinciannya diserahkan kepada pembaca.
\end{bukti}

% segment-id: o014.li2.chapter4.derived-functors-general.p008
Selanjutnya, misalkan
$Q_i:\cate{K}(\mathcal{A}_i)\to\cate{D}(\mathcal{A}_i)$ dan
$Q:\cate{K}(\mathcal{B})\to\cate{D}(\mathcal{B})$ merupakan funktor
lokalisasi ($i=1,2$). Tempatkan semua data ini dalam kerangka
Definisi~\ref{def:bifunctor-triangulated-localization}, dengan
$\mathcal{D}_1$, $\mathcal{D}_2$, dan $\mathcal{D}'$ berturut-turut
$\cate{K}^\star(\mathcal{A}_1)$, $\cate{K}^\star(\mathcal{A}_2)$, dan
$\cate{K}^\star(\mathcal{B})$, sedangkan kompleks-kompleks asiklik membentuk
$\mathcal{N}_1^\star$, $\mathcal{N}_2^\star$, dan $(\mathcal{N}')^\star$.

% segment-id: o014.li2.chapter4.derived-functors-general.d002
\begin{definisi}[Bifunktor turunan]\label{def:derived-bifunctor}
	\index{bifunktor turunan}
	\index[sym1]{RF} \index[sym1]{LF}
	Tetapkan $\star\in\{+,-,\hspace{0.8em}\}$ dan tinjau bifunktor aditif
	$F:\mathcal{A}_1\times\mathcal{A}_2\to\mathcal{B}$.
	\begin{itemize}
		\item Misalkan $\mathcal{B}$ mempunyai produk terhitung. Jika bifunktor
		turunan kanan dari $\cate{K}_{\Pi}F$ dalam pengertian
		Definisi~\ref{def:bifunctor-triangulated-localization},
		\[\cate{D}^\star(\mathcal{A}_1)\times\cate{D}^\star(\mathcal{A}_2)
		\to\cate{D}^\star(\mathcal{B}),\]
		ada, maka bifunktor tersebut disebut \emph{bifunktor turunan kanan} dari
		$F$ dan ditulis ${}^\star\mathrm{R}F$.
		\item Misalkan $\mathcal{B}$ mempunyai koproduk terhitung. Jika bifunktor
		turunan kiri dari $\cate{K}_{\oplus}F$,
		\[\cate{D}^\star(\mathcal{A}_1)\times\cate{D}^\star(\mathcal{A}_2)
		\to\cate{D}^\star(\mathcal{B}),\]
		ada, maka bifunktor tersebut disebut \emph{bifunktor turunan kiri} dari
		$F$ dan ditulis ${}^\star\mathrm{L}F$.
	\end{itemize}
	Tetap dituliskan
	${}^\star\mathrm{R}^nF:=\Hm^n\circ{}^\star\mathrm{R}F$ dan
	${}^\star\mathrm{L}_nF:=\Hm^{-n}\circ{}^\star\mathrm{L}F$.
\end{definisi}

% segment-id: o014.li2.chapter4.derived-functors-general.p009
Tetapkan $\star\in\{+,-,\hspace{0.8em}\}$ dan misalkan $\mathcal{I}_i$
subkategori bertriangulasi dari $\cate{K}^\star(\mathcal{A}_i)$, $i=1,2$.
Seperti biasa, jika $(\mathcal{I}_1,\mathcal{I}_2)$ merupakan pasangan
subkategori $\cate{K}^\star F$-injektif (atau projektif) dalam pengertian
Definisi~\ref{def:F-injective-bifunctor}, maka pasangan ini disebut
$F$-injektif (atau $F$-projektif). Proposisi~\ref{prop:localization-injective-bifunctor}
langsung memberikan kesimpulan berikut, dan semua isomorfismenya kanonik.
\index{subkategori injektif-$\cdots$} \index{subkategori projektif-$\cdots$}

% segment-id: o014.li2.chapter4.derived-functors-general.p010
\begin{itemize}
	\item Jika $(\mathcal{I}_1,\mathcal{I}_2)$ bersifat $F$-injektif, maka
	${}^\star\mathrm{R}F$ ada dan
	\begin{equation}\label{eqn:derived-bifunctor-1}
		(X_1, X_2) \in \Obj(\mathcal{I}_1) \times \Obj(\mathcal{I}_2) \implies {}^\star \mathrm{R}F(Q_1 X_1, Q_2 X_2) \simeq Q \cate{K}_{\Pi}F(X_1, X_2).
	\end{equation}
	\item Jika $(\mathcal{P}_1,\mathcal{P}_2)$ bersifat $F$-projektif, maka
	${}^*\mathrm{L}F$ ada dan
	\begin{equation}\label{eqn:derived-bifunctor-2}
		(X_1, X_2) \in \Obj(\mathcal{P}_1) \times \Obj(\mathcal{P}_2) \implies {}^\star \mathrm{L}F(Q_1 X_1, Q_2 X_2) \simeq Q \cate{K}_{\oplus}F(X_1, X_2).
	\end{equation}
\end{itemize}

% segment-id: o014.li2.chapter4.derived-functors-general.p011
Selanjutnya kita membahas suatu perkalian pada bifunktor turunan. Kita mulai
dengan versinya pada tingkat kompleks. Notasi tetap seperti di atas, tetapi
$F$ sekarang diasumsikan eksak kanan dalam setiap variabel. Untuk setiap
kompleks $X$, tuliskan $Z^n(X):=\Ker(d_X^n)$ dan
	$B^n(X):=\Image(d_X^{n-1})$. Untuk setiap $(p,q)\in\Z^2$ terdapat morfisme
	alami
\[ F\left( Z^p(X_1), Z^q(X_2) \right) \to Z^{p+q}\left( \cate{C}_{\oplus}F(X_1, X_2) \right). \]
Morfisme ini memetakan citra $F(B^p(X_1),Z^q(X_2))$ dan
$F(Z^p(X_1),B^q(X_2))$ ke
$B^{p+q}(\cate{C}_{\oplus}F(X_1,X_2))$. Dengan mengambil hasil bagi dan
menggunakan keeksakan kanan $F$, diperoleh morfisme kanonik
\[ \kappa: F\left(\Hm^p(X_1), \Hm^q(X_2)\right) \to \Hm^{p+q}\left( \cate{C}_{\oplus} F(X_1, X_2) \right). \]
Pembahasan sebelum Teorema Künneth untuk homologi
\ref{prop:Kunneth-homology} adalah kasus khusus konstruksi ini untuk
$F=\otimes_R$.

% segment-id: o014.li2.chapter4.derived-functors-general.p012
Sekarang ubah syaratnya: $F$ eksak kiri dalam setiap variabel. Dengan
mendualkan konstruksi di atas, diperoleh morfisme kanonik
\[ \lambda: \Hm^{p+q}\left( \cate{C}_{\Pi} F(X_1, X_2) \right) \to F\left( \Hm^p(X_1), \Hm^q(X_2) \right). \]
Keberadaan $\oplus$ atau $\prod$ terhitung yang dibahas tetap diasumsikan.
Karena hanya kohomologi yang diperhatikan di sini, $\cate{C}$ juga dapat
diganti dengan $\cate{K}$.

% segment-id: o014.li2.chapter4.derived-functors-general.pr002
\begin{proposisi}\label{prop:bifunctor-cup}
	Misalkan $F$ eksak kanan (atau eksak kiri) dalam setiap variabel,
	$\star\in\{+,-,\hspace{0.8em}\}$, dan terdapat pasangan $F$-projektif
	$(\mathcal{P}_1,\mathcal{P}_2)$ (atau pasangan $F$-injektif
	$(\mathcal{I}_1,\mathcal{I}_2)$). Maka terdapat morfisme kanonik
	\[ F\left(\Hm^p(\cdot), \Hm^q(\cdot) \right) \to \Hm^{p+q} {}^\star \mathrm{L}F \quad \text{atau} \quad \Hm^{p+q} {}^\star \mathrm{R}F \to F\left(\Hm^p(\cdot), \Hm^q(\cdot) \right), \]
	dengan kedua ruas merupakan funktor dari
	$\cate{D}^\star(\mathcal{A}_1)\times\cate{D}^\star(\mathcal{A}_2)$ ke
	$\mathcal{B}$. Morfisme-morfisme tersebut dicirikan oleh diagram komutatif
% segment-id: o014.li2.chapter4.derived-functors-general.g005
	\begin{equation*}
		\begin{tikzcd}
			& \Hm^{p+q} \cate{K}_{\oplus} F(X_1, X_2) \\
			F\left(\Hm^p(X_1), \Hm^q(X_2) \right) \arrow[r] \arrow[ru, "\kappa"] & \Hm^{p+q} {}^\star \mathrm{L}F(Q_1 X_1, Q_2 X_2) \arrow[u, "\mathrm{can}"']
		\end{tikzcd}
	\end{equation*}
	atau
% segment-id: o014.li2.chapter4.derived-functors-general.g006
	\begin{equation*}
		\begin{tikzcd}[column sep=small]
			& \Hm^{p+q} \cate{K}_{\Pi} F(X_1, X_2) \arrow[d, "\mathrm{can}"] \arrow[ld, "\lambda"'] \\
			F\left(\Hm^p(X_1), \Hm^q(X_2) \right) & \Hm^{p+q} {}^\star \mathrm{R}F(Q_1 X_1, Q_2 X_2), \arrow[l]
		\end{tikzcd}
	\end{equation*}
	dengan $X_i\in\Obj(\cate{K}^\star(\mathcal{A}_i))$. Panah yang diberi
	label $\mathrm{can}$ adalah morfisme kanonik yang menyertai funktor turunan
	sebagai ekstensi Kan (setelah menerapkan $\Hm^{p+q}$); lihat
	\eqref{eqn:L-R-can} untuk kasus satu variabel.
\end{proposisi}
% segment-id: o014.li2.chapter4.derived-functors-general.pf004
\begin{bukti}
	Berdasarkan dualitas, cukup dibahas kasus yang melibatkan $\kappa$. Perhatikan
	bahwa, selain $\cate{K}_{\oplus}F(X_1,X_2)$, dua simpul lain dalam diagram
	hanya bergantung pada $Q_1X_1$ dan $Q_2X_2$.

	Tetapkan pasangan $(\mathcal{P}_1,\mathcal{P}_2)$. Karena semua morfisme dalam diagram
	bersifat funktorial dan
	$\mathcal{P}_i/(\mathcal{P}_i\cap\mathcal{N}^\star_i)$ ekuivalen dengan
	$\cate{D}^\star(\mathcal{A}_i)$, cukup menentukan morfisme yang dicari untuk
	$X_i\in\Obj(\mathcal{P}_i)$; hal ini dapat diperiksa dengan menggambar
	diagram. Namun, dalam kasus ini
	${}^\star\mathrm{L}F(Q_1X_1,Q_2X_2)=\cate{K}_{\oplus}F(X_1,X_2)$, sehingga
	pernyataannya langsung.
\end{bukti}

% segment-id: o014.li2.chapter4.derived-functors-general.p013
Singkatnya, semuanya dapat dihitung dengan mengambil resolusi, dan hasil
akhirnya tidak bergantung pada pilihan resolusi.
